\documentclass[border=10pt]{ctexart} % 使用ctexart文档类
\usepackage{tikz}
\usepackage{circuitikz}

\begin{document}
\pagestyle{empty} % 去除页码

\begin{circuitikz}[
    american,
    thick,
    scale=1.0,
    every label/.style={font=\footnotesize}
]

% -------------------------- 1. 核心电路区（仅保留编号，无多余文字） --------------------------
% 参考节点：定压外输终端
\node[ground] (GND) at (0,0) {};

% 管汇节点
\node[circ] (nodeA) at (0,4) {};
\node[circ] (nodeB) at (5,4) {};

% 管线（电阻）
\draw (GND) to[R, l_=$R_1$] (nodeA);
\draw (GND) to[R, l^=$R_2$] (nodeB);
\draw (nodeA) to[R, l^=$R_3$] (nodeB);

% 生产井（电流源）
\draw (nodeA) to[I, l_=$Q_1$, invert] (0,6) node[ground] {};
\draw (nodeB) to[I, l^=$Q_2$, invert] (5,6) node[ground] {};

% -------------------------- 2. 嵌入图内的图例（安全绘制，无嵌套冲突） --------------------------
% 图例标题
\node[font=\bfseries\footnotesize, anchor=west] at (-1.5, -1.2) {图 例：};

% 定义图例起始坐标
\coordinate (leg) at (-1.5, -1.8);

% --- 第一列：元件符号说明 ---
% R1
\node[anchor=west, font=\footnotesize] at (leg) {$R_1$：};
\node[anchor=west, font=\footnotesize] at ([xshift=3.2em]leg) {A区块外输总管};

% R2
\node[anchor=west, font=\footnotesize] at ([yshift=-1em]leg) {$R_2$：};
\node[anchor=west, font=\footnotesize] at ([xshift=3.2em, yshift=-1em]leg) {B区块外输总管};

% R3
\node[anchor=west, font=\footnotesize] at ([yshift=-2em]leg) {$R_3$：};
\node[anchor=west, font=\footnotesize] at ([xshift=3.2em, yshift=-2em]leg) {区块间互联管线};

% Q1/Q2
\node[anchor=west, font=\footnotesize] at ([yshift=-3em]leg) {$Q_1, Q_2$：};
\node[anchor=west, font=\footnotesize] at ([xshift=3.2em, yshift=-3em]leg) {区块生产井组};

% --- 第二列：核心类比关系 ---
\coordinate (leg2) at (3, -1.8);

\node[anchor=west, font=\footnotesize] at (leg2) {电压 $U$ $\leftrightarrow$ 压力 $\Delta P$};
\node[anchor=west, font=\footnotesize] at ([yshift=-1em]leg2) {电流 $I$ $\leftrightarrow$ 流量 $Q$};
\node[anchor=west, font=\footnotesize] at ([yshift=-2em]leg2) {电阻 $R$ $\leftrightarrow$ 流动阻力 $R_{\mathrm{flow}}$};
\node[anchor=west, font=\footnotesize] at ([yshift=-3em]leg2) {接地 $\leftrightarrow$ 定压外输终端};

\end{circuitikz}

\end{document}